% Augmented Lagrangian for hybrid PINN training — advisor slides
% Compile:  pdflatex augmented_lagrangian.tex   (run twice)
%       or: open in Overleaf
\documentclass[aspectratio=169,11pt]{beamer}
\usetheme{Madrid}
\usecolortheme{seahorse}
\setbeamertemplate{navigation symbols}{}
\setbeamertemplate{footline}[frame number]

\usepackage{amsmath,amssymb}
\usepackage{booktabs}

\newcommand{\Ldata}{\mathcal{L}_{\text{data}}}
\newcommand{\vhat}{\hat{v}}

\title[Augmented Lagrangian for PINNs]{Why an Augmented Lagrangian for our PINN training}
\subtitle{Fixing loss-balancing drift at the source}
\author{PINN Option-Pricing project}
\date{\today}

\begin{document}

\begin{frame}
  \titlepage
\end{frame}

% ---------------------------------------------------------------
\begin{frame}{1. The problem: our loss is a ``penalty method''}
  Today we minimise a \textbf{weighted sum} with auto-tuned weights (grad-norm):
  \[
    \mathcal{L} \;=\; \lambda_{\text{pde}}\mathcal{L}_{\text{pde}}
      + \lambda_{\text{tc}}\mathcal{L}_{\text{tc}}
      + \lambda_{\text{bc}}\mathcal{L}_{\text{bc}}
      + \lambda_{\text{data}}\Ldata .
  \]
  In optimisation terms this is the \alert{penalty method} for a constrained problem.

  \begin{block}{Its built-in flaw}
    To truly \emph{enforce} a term you must send its weight $\to\infty$; large
    weights make the optimisation stiff and unstable.
  \end{block}

  \begin{exampleblock}{What we observe}
    $\lambda_{\text{data}}$ runs away to ${\sim}10^{5}$, training drifts:
    physics-only RMSE $\$8.76 \rightarrow$ hybrid $\$10$+, and no-arbitrage
    violations explode. The weights never settle.
  \end{exampleblock}
\end{frame}

% ---------------------------------------------------------------
\begin{frame}{2. Reframe: it is really a \emph{constrained} problem}
  We don't want a blend of physics and data. We want:
  \begin{quote}
    \textbf{minimise the market-data misfit \emph{subject to} obeying the
    Black--Scholes PDE, the boundary/terminal conditions, and no-arbitrage.}
  \end{quote}
  \[
  \begin{aligned}
    \min_{\theta}\quad & \Ldata(\theta) &&\text{\small(objective: fit the market)}\\
    \text{s.t.}\quad
      & \mathcal{L}_{\text{pde}}(\theta)=0,\;
        \mathcal{L}_{\text{tc}}(\theta)=0,\;
        \mathcal{L}_{\text{bc}}(\theta)=0 &&\text{\small(equality constraints)}\\
      & \underbrace{\mathrm{ReLU}(-\partial^2_{mm}\vhat)}_{\text{butterfly}}\le 0,\;
        \underbrace{\mathrm{ReLU}(-\partial_{\tau}\vhat)}_{\text{calendar}}\le 0
        &&\text{\small(no-arbitrage)}
  \end{aligned}
  \]
  \begin{block}{The key shift}
    Data is the \textbf{objective}; physics and no-arbitrage are \textbf{constraints}.
    The weighted sum throws this structure away by treating them as co-equal losses.
  \end{block}
\end{frame}

% ---------------------------------------------------------------
\begin{frame}{3. What the Augmented Lagrangian does}
  Introduce a \alert{multiplier} $\mu_k$ per constraint $c_k$ and add a penalty:
  \[
    \mathcal{L}_{A}(\theta,\mu;\rho)=\Ldata(\theta)
      \;+\;\sum_k \mu_k\,c_k(\theta)
      \;+\;\tfrac{\rho}{2}\sum_k c_k(\theta)^2 .
  \]
  Alternate two cheap steps:
  \begin{itemize}
    \item descend on $\theta$ (a few gradient steps);
    \item update multipliers \;$\mu_k \leftarrow \mu_k + \rho\,c_k$\;
          \big(inequalities: $\mu_k \leftarrow \max(0,\ \mu_k+\rho\,c_k)$\big).
  \end{itemize}

  \begin{block}{Why it cures the drift}
    The multiplier \textbf{accumulates the violation} and does the job that
    $\rho\to\infty$ would do --- but with a \textbf{fixed, moderate} $\rho$.
    It is \alert{self-limiting}: once $c_k\!\approx\!0$ the update adds ${\approx}0$,
    so $\mu_k$ stops growing. No runaway, by construction.
  \end{block}

  \centering\small
  \emph{Analogy:} penalty method $=$ ever-larger fines until compliance
  ($\to\infty$); \quad augmented Lagrangian $=$ a \textbf{thermostat} that settles
  at exactly the pressure needed, then holds.
\end{frame}

% ---------------------------------------------------------------
\begin{frame}{4. What it enables for us}
  \begin{columns}[T]
    \begin{column}{0.54\textwidth}
      \begin{itemize}
        \item \textbf{Fixes drift at the source} $\Rightarrow$ trajectory
              converges $\Rightarrow$ final-epoch reliable $\Rightarrow$
              \emph{we can drop the validation set} and grow the training window.
        \item \textbf{Encodes the right structure}: physics \& no-arbitrage are
              constraints, data is the objective.
        \item \textbf{Unifies no-arbitrage}: just another inequality constraint,
              active only where violated --- self-modulating.
        \item \textbf{Scales to Stage 2}: $\sigma$-regularisation and the second
              network are simply more constraints.
      \end{itemize}
    \end{column}
    \begin{column}{0.46\textwidth}
      \footnotesize
      \begin{tabular}{@{}ll@{}}
        \toprule
        & uses \\
        \midrule
        current grad-norm & per-term gradient \emph{norms} \\
        augmented Lagrangian & constraint \emph{values} \\
        \bottomrule
      \end{tabular}
      \vspace{0.4em}
      \begin{block}{\footnotesize Cost}
        \footnotesize One robust knob $\rho$. Multiplier updates use constraint
        \emph{values}, not per-term gradients --- so it is \textbf{cheaper and more
        stable} than what we run now.
      \end{block}
      \scriptsize PINN precedent: PECANN (Basir \& Senocak, 2022); AL-PINN
      (Son et al., 2022). Method: Hestenes \& Powell, 1969.
    \end{column}
  \end{columns}
\end{frame}

\end{document}
